% Historical working note: verify statements against the current Rocq sources.
\documentclass[11pt]{article}

\usepackage[T1]{fontenc}
\usepackage[utf8]{inputenc}
\usepackage{amsmath,amssymb,amsthm}

\newtheorem{theorem}{Tétel}
\newtheorem{definition}{Definíció}
\newtheorem{proposition}{Propozíció}
\newtheorem{remark}{Megjegyzés}

\newenvironment{proofsketch}
  {\par\noindent\textit{Bizonyításvázlat.}}
  {\hfill$\square$\par}
\newenvironment{coqblock}
  {\begin{quote}\small\ttfamily\obeylines}
  {\end{quote}}

\newcommand{\SetoidCat}{\mathbf{SetoidCat}}
\newcommand{\EpsMod}{\mathbf{EpsMod}_{\Sigma}}
\newcommand{\Raw}{\mathbf{Raw}_{\Sigma}}
\newcommand{\Step}{\mathsf{Step}}
\newcommand{\TM}{\mathbb{T}}
\newcommand{\Hom}{\operatorname{Hom}}
\newcommand{\Obj}{\operatorname{Obj}}
\newcommand{\Ty}{\operatorname{Ty}}
\newcommand{\Ctx}{\operatorname{Ctx}}
\newcommand{\BT}{\operatorname{BT}}
\newcommand{\HD}{\operatorname{HD}}
\newcommand{\WF}{\operatorname{WF}}
\newcommand{\AWF}{\operatorname{AWF}}
\newcommand{\Fib}{\operatorname{Fib}}
\newcommand{\sem}[1]{[\![#1]\!]}
\newcommand{\id}{\operatorname{id}}
\newcommand{\comp}{\mathbin{\circ}}
\newcommand{\eps}{\varepsilon}

\title{Inicialitás és fixponttételek az EpsilonFull termmodellre}
\author{}
\date{}

\begin{document}
\maketitle

\section{Rögzített szignatúra}

Legyen rögzítve egy EpsilonFull szignatúra
\[
  \Sigma=(\mathsf{Sign},\mathsf{Func},\mathsf{Pred},
          \mathsf{HFunc},\mathsf{HPred},\mathsf{SM}),
\]
ahol $\mathsf{HFunc}$ a függvényjelek, $\mathsf{HPred}$ a
predikátumjelek, $\mathsf{SM}$ pedig a Skolem--epsilon választások
szignatúraadata. Az alábbi állítások erre a rögzített szignatúrára
vonatkoznak.

\section{Az EpsilonFull-modellek kategóriája}

\begin{definition}[EpsilonFull-modell]
Egy EpsilonFull-modell $M$ a formalizációban a következő adatokból áll:
\[
\begin{gathered}
  \mathcal C_M \text{ setoid kategória}, \qquad
  \mathcal C_M \text{ karteziánus struktúrával},                                    \\
  \text{kétsortú bázisinterpretáció},\quad
  \text{koherens bázismodell},                                                       \\
  p_M:\mathcal E_M\to \mathcal C_M
  \text{ jó és strict fibráció},                                                     \\
  \text{hiperlogikai interpretáció},\quad
  \text{AdjointEpsilonFullFibration struktúra}.
\end{gathered}
\]
Coq-ban ez a \texttt{EpsilonFullModel} osztály.
\end{definition}

Ha $I$ bázistípus, akkor $\sem{I}_M$ jelöli az $I$ interpretációját
$\mathcal C_M$-ben. Ha $\varphi$ formula $I$ felett, akkor
$\sem{\varphi}_M\in \Fib_M(I)$, ahol
\[
  \Fib_M(I)=\mathcal E_M{}_{\sem{I}_M}
\]
a megfelelő rostkategória.

\begin{definition}[Modellmorfizmus]
Legyen $M,N$ két EpsilonFull-modell. Egy
$F:M\to N$ morfizmus három pontonként definiált térképcsaládból áll:
\[
\begin{aligned}
F_b^{I,A} &: \mathcal C_M(\sem I_M,\sem A_M)
          \longrightarrow
          \mathcal C_N(\sem I_N,\sem A_N),\\
F_f^{I,\varphi,\psi} &: \Fib_M(I)(\sem\varphi_M,\sem\psi_M)
          \longrightarrow
          \Fib_N(I)(\sem\varphi_N,\sem\psi_N),\\
F_s^{I,\Gamma,\varphi} &: \Fib_M(I)(\sem\Gamma_M,\sem\varphi_M)
          \longrightarrow
          \Fib_N(I)(\sem\Gamma_N,\sem\varphi_N).
\end{aligned}
\]
Ezek proper térképek a megfelelő setoid-hom ekvivalenciákra nézve. A
fiber-komponens megőrzi az identitást és a kompozíciót, a bázis- és
hiperoldalon pedig teljesíti a soundness feltételeket:
\[
  F_b^{I,A}(\sem{t}_{M,d})\simeq \sem{t}_{N,d}
  \qquad
  (d:\BT(I,t,A)),
\]
illetve
\[
  F_s^{I,\Gamma,\varphi}(\sem{p}_{M,d})
  \simeq
  \sem{p}_{N,d}
  \qquad
  (d:\HD(I,\Gamma,p,\varphi)).
\]
\end{definition}

\begin{definition}[Morfizmusegyenlőség]
Két morfizmus $F,G:M\to N$ ekvivalens, jelölésben $F\sim G$, ha
mindhárom komponensen pontonként egyenlők a célbeli hom-ekvivalencia
szerint:
\[
\begin{aligned}
&\forall I,A,f,\quad
  F_b^{I,A}(f)\simeq G_b^{I,A}(f),\\
&\forall I,\varphi,\psi,f,\quad
  F_f^{I,\varphi,\psi}(f)\simeq G_f^{I,\varphi,\psi}(f),\\
&\forall I,\Gamma,\varphi,f,\quad
  F_s^{I,\Gamma,\varphi}(f)\simeq G_s^{I,\Gamma,\varphi}(f).
\end{aligned}
\]
Ez a Coq-beli \texttt{EpsilonFullModelHom\_equiv}. Ez nem triviális
reláció.
\end{definition}

\begin{proposition}
Az EpsilonFull-modellek, a fenti morfizmusok és a fenti
morfizmusegyenlőség large setoid kategóriát alkotnak:
\[
  \EpsMod.
\]
Coq-ban:
\begin{coqblock}
EpsilonFullModel\_LargeSetoidCategory
\end{coqblock}
\end{proposition}

\section{Inicialitás}

Jelölje $\TM$ a termmodellt
(\texttt{TermModel\_EpsilonFullModel}).

\begin{theorem}[A termmodell inicialitása]
A termmodell iniciális objektum az EpsilonFull-modellek large setoid
kategóriájában:
\[
  \TM \text{ iniciális objektum } \EpsMod\text{-ban}.
\]
Formálisan, minden $M\in\Obj(\EpsMod)$ modellre létezik morfizmus
\[
  \eta_M:\TM\longrightarrow M,
\]
és ha $F,G:\TM\to M$, akkor
\[
  F\sim G.
\]
Coq-ban:
\begin{coqblock}
TermModel\_to M : EpsilonFullModelHom TermModel M
TermModel\_LargeInitial : LargeInitial TermModel
TermModel\_initial\_unique :
  forall M F G, EpsilonFullModelHom\_equiv F G
\end{coqblock}
\end{theorem}

\begin{proofsketch}
Először megadjuk a kanonikus morfizmust $\eta_M:\TM\to M$.
Bázison egy termmodellbeli nyilat visszaolvasunk egy jól típusozott
bázistermmé, majd ezt interpretáljuk $M$-ben:
\[
  (\eta_M)_b(f)
  =
  \sem{\operatorname{deriv}(\operatorname{read}_b(f))}_M.
\]
A fiber- és sequent-komponensek hasonlóan keletkeznek: a termmodellbeli
homot zárt hiperlevezetési termként olvassuk vissza, majd az EpsilonFull
soundness adja az $M$-beli homot.

Az egyediséghez legyen $F:\TM\to M$ tetszőleges morfizmus. A bázis
komponensen a formalizáció bizonyítja
\[
  F_b^{I,A}(f)\simeq (\eta_M)_b^{I,A}(f).
\]
Ennek két összetevője van. Egyrészt minden termmodellbeli bázishom
ekvivalens a visszaolvasott derivációjának termmodellbeli
interpretációjával:
\[
  f\simeq \sem{\operatorname{deriv}(\operatorname{read}_b(f))}_{\TM}.
\]
Másrészt $F$ bázis-soundness klauzulája a visszaolvasott deriváció
interpretációját $M$-beli interpretációba viszi.

A fiber- és sequent-komponenseken az egyediség a jó fibráció részét képező
thinness feltételből következik:
\[
  \forall I,\quad \Fib_M(I)\text{ thin}.
\]
Ezért azonos forrás és cél között bármely két rostbeli hom ekvivalens.
Összegezve a három pontonkénti komponensekvivalenciát kapjuk, vagyis
$F\sim G$.
\end{proofsketch}

\section{A nyers hordozók kategóriája és a szintaxisgenerátor}

\begin{definition}[Nyers hordozó]
Egy nyers EpsilonFull-hordozó egy pár
\[
  X=(X_b,X_h),
\]
ahol
\[
  X_b(I,t,A):\mathsf{Type},
  \qquad
  X_h(I,\Gamma,p,\varphi):\mathsf{Type}.
\]
Itt $X_b(I,t,A)$ az $I\vdash t:A$ alakú bázisadatok, míg
$X_h(I,\Gamma,p,\varphi)$ az
$I;\Gamma\vdash p:\varphi$ alakú hiperadatok hordozója.
\end{definition}

\begin{definition}[Nyers hordozók morfizmusa]
Egy $u:X\to Y$ morfizmus indexet tartó függvénypár:
\[
\begin{aligned}
u_b^{I,t,A} &: X_b(I,t,A)\to Y_b(I,t,A),\\
u_h^{I,\Gamma,p,\varphi} &: X_h(I,\Gamma,p,\varphi)\to
                           Y_h(I,\Gamma,p,\varphi).
\end{aligned}
\]
Két ilyen morfizmus akkor egyenlő, ha ezek a függvények minden indexen
pontonként proposicionálisan egyenlők.
\end{definition}

\begin{proposition}
A nyers hordozók large setoid kategóriát alkotnak:
\[
  \Raw.
\]
Coq-ban:
\begin{coqblock}
EpsilonFullRawCarrier\_LargeSetoidCategory
\end{coqblock}
\end{proposition}

\begin{definition}[Az egy-lépéses szintaxisgenerátor]
Definiáljuk a
\[
  \Step:\Raw\longrightarrow\Raw
\]
endofunktort objektumokon a következőképpen:
\[
\begin{aligned}
  (\Step X)_b(I,t,A) &= \mathsf{BaseStep}_X(I,t,A),\\
  (\Step X)_h(I,\Gamma,p,\varphi) &= \mathsf{HyperStep}_X(I,\Gamma,p,\varphi).
\end{aligned}
\]
A $\mathsf{BaseStep}_X$ családot a bázisszabályok egyetlen külső
alkalmazásai generálják:
\[
\begin{array}{rcl}
\mathsf{id}      &:& \mathsf{BaseStep}_X(A,\mathsf{id},A),\\
\mathsf{comp}    &:& X_b(B,g,C)\times X_b(A,f,B)
                  \to \mathsf{BaseStep}_X(A,g\circ f,C),\\
\mathsf{unit}    &:& \mathsf{BaseStep}_X(A,\mathsf{unit},1),\\
\mathsf{pair}    &:& X_b(A,f,B)\times X_b(A,g,C)
                  \to \mathsf{BaseStep}_X(A,\langle f,g\rangle,B\times C),\\
\mathsf{fst}     &:& \mathsf{BaseStep}_X(A\times B,\pi_1,A),\\
\mathsf{snd}     &:& \mathsf{BaseStep}_X(A\times B,\pi_2,B),\\
\mathsf{const}   &:& \mathsf{BaseStep}_X(I,\mathsf{const}_A(s),A),\\
\mathsf{func}    &:& \mathsf{BaseStep}_X(\operatorname{dom}(f),
                         \mathsf{func}(f),\operatorname{cod}(f)).
\end{array}
\]
A $\mathsf{HyperStep}_X$ családot a hiperlevezetési szabályok egyetlen
külső alkalmazásai generálják. A rekurzív hiper-premisszák helyén $X_h$,
a bázispremisszák helyén $X_b$ szerepel. Például:
\[
\begin{aligned}
\mathsf{app} &: X_h(I,\Gamma,p,\varphi\Rightarrow\psi)
               \times X_h(I,\Gamma,q,\varphi)
               \to X_h'(I,\Gamma,p\,q,\psi),\\
\forall\mathsf{E} &: X_h(I,\Gamma,p,\forall_A\varphi)
               \times X_b(I,f,A)
               \to X_h'(I,\Gamma,p[f],\varphi[(\id,f)]),\\
\exists\mathsf{I} &: X_b(I,f,A)
               \times X_h(I,\Gamma,p,\varphi[(\id,f)])
               \to X_h'(I,\Gamma,\exists\mathsf{I}(f,p),\exists_A\varphi),\\
\eps\mathsf{I} &: X_b(I,f,A)
               \times X_h(I,\Gamma,p,\varphi[(\id,f)])
               \to
               X_h'(I,\Gamma,\eps\mathsf{I}(f,p),
                    \varphi[(\id,\eps_{A,\varphi})]).
\end{aligned}
\]
Itt $X_h'$ rövidítés a $(\Step X)_h$ családra. A többi konstruktor
pontosan a \texttt{h\_derivation} megfelelő szabályának egy-lépéses
változata.
\end{definition}

Morfizmusokon $\Step(u)$ úgy hat, hogy az egy-lépéses konstruktorok
azonnali rekurzív premisszáiban alkalmazza $u_b$-t, illetve $u_h$-t.
Így $\Step$ megőrzi az identitást és a kompozíciót. Coq-ban:
\begin{coqblock}
EpsilonFullRawOneStep\_Endofunctor :
  LargeEndofunctor EpsilonFullRawCarrier\_LargeSetoidCategory
\end{coqblock}

\section{A nyers szintaxis fixponttétele}

Definiáljuk a term hordozót $T\in\Raw$:
\[
  T_b(I,t,A)=\BT(I,t,A),
  \qquad
  T_h(I,\Gamma,p,\varphi)=\HD(I,\Gamma,p,\varphi).
\]
Coq-ban ez a \texttt{TermModelRawCarrier}.

\begin{theorem}[A nyers term hordozó fixpont]
A term hordozó izomorf az egy-lépéses szintaxisgenerátor szerinti képével:
\[
  \Step(T)\cong T
  \qquad\text{a }\Raw\text{ large setoid kategóriában.}
\]
Coq-ban:
\begin{coqblock}
TermModelRawCarrier\_OneStepIsomorphism
\end{coqblock}
\end{theorem}

\begin{proofsketch}
Az izomorfizmus két iránya
\[
  \rho:\Step(T)\to T,
  \qquad
  \upsilon:T\to \Step(T).
\]
A $\rho$ morfizmus a \texttt{RawBaseOneStep} konstruktorokat a megfelelő
\texttt{base\_typing} konstruktorokra, a \texttt{RawHyperOneStep}
konstruktorokat pedig a megfelelő \texttt{h\_derivation} konstruktorokra
küldi.

Az $\upsilon$ morfizmus egy adott \texttt{base\_typing}, illetve
\texttt{h\_derivation} derivációt a külső konstruktorán esetszétválasztva
egy one-step objektummá bont.

Az inverztörvények:
\[
  \rho\circ\upsilon \sim \id_T,
  \qquad
  \upsilon\circ\rho \sim \id_{\Step(T)}.
\]
Mindkét egyenlőség a $\Raw$ kategória pontonkénti
morfizmusegyenlősége szerint értendő. A bizonyítás konstruktor szerinti
esetszétválasztás után definicionális redukció.
\end{proofsketch}

\begin{remark}
Ez Lambek-típusú állítás egy konkrét endofunktorra. A formalizáció itt nem
állít általános Lambek-lemmát tetszőleges kategóriára és tetszőleges
endofunktorra.
\end{remark}

\section{A termmodell hom-szintű fixponttétele}

A nyers fixponttétel mellett a formalizáció a termmodell tényleges homjaira
is kimondja a roll--unroll tulajdonságot.

\subsection{Bázis}

Legyen
\[
  \Hom_b^{\TM}(I,A)
  =
  \mathcal C_{\TM}(\sem I_{\TM},\sem A_{\TM}).
\]
Van egy one-step típus $\mathsf{BaseOneStepHom}(I,A)$, valamint térképek
\[
  r_b:\mathsf{BaseOneStepHom}(I,A)\to \Hom_b^{\TM}(I,A),
  \qquad
  u_b:\Hom_b^{\TM}(I,A)\to \mathsf{BaseOneStepHom}(I,A).
\]

\begin{theorem}[Szemantikus bázis fixpont]
Minden $f\in \Hom_b^{\TM}(I,A)$ esetén
\[
  r_b(u_b(f))\simeq f.
\]
Továbbá minden $s\in \mathsf{BaseOneStepHom}(I,A)$ esetén
\[
  u_b(r_b(s))\approx s,
\]
ahol az utóbbi ekvivalencia a roll által indukált one-step ekvivalencia.
Coq-ban:
\begin{coqblock}
TermModel\_semantic\_base\_roll\_unroll
TermModel\_semantic\_base\_unroll\_roll
\end{coqblock}
\end{theorem}

\begin{proofsketch}
Egy termmodellbeli bázishomot a \texttt{term\_model\_base\_hom}
visszaolvas jól típusozott bázistermmé. A
$r_b(u_b(f))\simeq f$ állítás a bázis soundness ekvacionális alakját és a
termmodell reinterpetációs lemmáját használja. A másik irány ebből következik
a one-step ekvivalencia definíciója szerint.
\end{proofsketch}

\subsection{Fiber és sequent}

Legyen
\[
  \Hom_f^{\TM}(I,\varphi,\psi)
  =
  \Fib_{\TM}(I)(\sem\varphi_{\TM},\sem\psi_{\TM})
\]
és
\[
  \Hom_s^{\TM}(I,\Gamma,\varphi)
  =
  \Fib_{\TM}(I)(\sem\Gamma_{\TM},\sem\varphi_{\TM}).
\]
A formalizáció definiál one-step típusokat
\[
  \mathsf{FiberOneStep}(I,\varphi,\psi),
  \qquad
  \mathsf{SequentOneStep}(I,\Gamma,\varphi),
\]
és a megfelelő roll--unroll térképeket.

\begin{theorem}[Szemantikus fiber és sequent fixpont]
Jólformáltsági feltételek mellett minden fiber-homra
\[
  r_f(u_f(f))\simeq f,
  \qquad
  u_f(r_f(s))\approx s,
\]
és minden sequent-homra
\[
  r_s(u_s(f))\simeq f,
  \qquad
  u_s(r_s(s))\approx s.
\]
Coq-ban:
\begin{coqblock}
TermModel\_semantic\_fiber\_roll\_unroll
TermModel\_semantic\_fiber\_unroll\_roll
TermModel\_semantic\_sequent\_roll\_unroll
TermModel\_semantic\_sequent\_unroll\_roll
\end{coqblock}
\end{theorem}

\begin{proofsketch}
A fiber- és sequent-homok one-step reprezentációja zárt hiperlevezetési
termekkel történik. Az egyenlőségi állítások a termmodell megfelelő
rostjainak thinness tulajdonságából következnek:
\[
  \Fib_{\TM}(I)\text{ thin}.
\]
Ez pontosan a \texttt{fiber\_is\_thin} feltétel termmodellre specializált
alakja.
\end{proofsketch}

\begin{proposition}[Összefoglaló Coq-osztály]
A három szemantikus komponens és a raw one-step fixpont együtt:
\begin{coqblock}
TermModel\_ConcreteOneStepFixedPoint :
  TermModelConcreteOneStepFixedPoint
\end{coqblock}
\end{proposition}

\section{A három eredmény viszonya}

Az inicialitás és a fixponttételek különböző kategóriákban élnek, de ugyanazt
a matematikai képet támasztják alá.

\begin{enumerate}
\item Az inicialitás a teljes EpsilonFull szemantikai modellek kategóriájában
áll:
\[
  \TM \text{ iniciális } \EpsMod\text{-ban}.
\]
Ez az interpretációk univerzális tulajdonsága.

\item A kategóriai fixpont a nyers, indexelt bizonyítás-hordozók
kategóriájában áll:
\[
  \Step(T)\cong T \quad \text{a }\Raw\text{-ban}.
\]
Ez a szintaxisgenerálás algebrai fixponttulajdonsága.

\item A szemantikus fixpont a termmodell tényleges homjaira vonatkozik:
a bázis-, fiber- és sequent-homok egy-lépéses szintaktikus alakra bonthatók,
majd visszagörgethetők a megfelelő hom-ekvivalencia szerint.
\end{enumerate}

Fontos negatív állítások:
\[
\begin{gathered}
\text{nem használunk triviális morfizmusegyenlőséget;}\\
\text{nem állítunk általános Lambek-lemmát;}\\
\text{a }\Step\text{ endofunktor nem a teljes }\EpsMod\text{ kategórián él,}\\
\text{hanem a nyers hordozók }\Raw\text{ kategóriáján.}
\end{gathered}
\]

\end{document}
